% Generated draft: requires mathematical and source review.
\documentclass[11pt]{article}
\usepackage[margin=25mm]{geometry}
\usepackage{fontspec}
\setmainfont{Latin Modern Roman}
\usepackage{amsmath,amssymb,mathtools,unicode-math}
\setmathfont{Latin Modern Math}
\usepackage{enumitem,needspace,xurl,hyperref}
\setlength{\parindent}{0pt}
\setlength{\parskip}{5pt}
\setlength{\emergencystretch}{3em}
\newcommand{\N}{\mathbb{N}}
\newcommand{\Z}{\mathbb{Z}}
\newcommand{\Q}{\mathbb{Q}}
\newcommand{\R}{\mathbb{R}}
\newcommand{\C}{\mathbb{C}}
\newcommand{\F}{\mathbb{F}}
\newcommand{\A}{\mathbb{A}}
\newcommand{\PP}{\mathbb{P}}
\newcommand{\E}{\mathbb{E}}
\newcommand{\eps}{\varepsilon}
\newcommand{\dd}{\,\mathrm{d}}
\begin{document}
\section*{UnsolvedMath ID: 30000768}
\textbf{Source problem code:} catalogue entry
\subsection{Definitions and mathematical statement}
Let \((K,\partial)\) be a characteristic-zero differential function field, with constant field \(C=\ker\partial\), in the intended complex setting \(C=\C\). The supplied record does not fix the base variety, number of derivations, or differential overfield. For a smooth commutative algebraic group \(G/K\), an algebraic \(D\)-group structure is a compatible algebraic group section \(s:G\to\tau_\partial G\) of the differential prolongation. If \((L,\partial)\supseteq(K,\partial)\) is a differential extension, the difference between the derivative of \(y\in G(L)\) and \(s(y)\), translated to the identity, gives the logarithmic derivative
\[
 \partial\log_G:G(L)\longrightarrow\operatorname{Lie}(G)(L).
\]
The induced \(D\)-structure on the Lie algebra gives a differential operator \(\partial_{\operatorname{Lie}(G)}\). Its horizontal vectors and the solutions of \(\partial\log_G(y)=0\) are the respective constant-solution objects.

In this language, the requested extension of Ax's generalized Schanuel theorem concerns transcendence bounds and differential-Galois descriptions for logarithmic and exponential equations, schematically
\[
 \partial\log_G(y)=a,
 \qquad
 \partial\log_G(y)=\partial_{\operatorname{Lie}(G)}(x).
\]
The problem asks to extend the constant commutative algebraic-group theory over \(\C\) to non-isoconstant algebraic \(D\)-groups over a function field, including both formulations and their differential-Galois interpretation. For a constant group \(G_0/C\), base change to \(K\) with its constant \(D\)-structure supplies the reference constant case. Under the differential-group meaning, isoconstant means becoming isomorphic to such a constant \(D\)-group after an allowed differential field extension; if the source instead uses isoconstancy of the underlying algebraic group, that distinction must be resolved.

This is a research-program request, not a fully specified quantified inequality in the supplied record. In particular, it supplies no numerical transcendence-degree bound, base field relative to which it is measured, allowed class of nonconstant groups, or independence condition modulo proper algebraic or \(D\)-subgroups. No claim that every solution of either displayed schematic equation is generic is intended.
\subsection{Short English statement}
The requested theory would extend the constant-group Ax transcendence framework to varying algebraic differential groups over function fields, with logarithmic, exponential, and differential-Galois versions. The record leaves the precise conjectural bounds and nondegeneracy hypotheses unstated.
\subsection{Sources}
\begin{itemize}
\item Local UnsolvedMath registry, permanent ID \(30000768\), exact source code \texttt{catalogue entry}, proposed year 2007; no individual proposer is supplied. The requested extension is attributed to the framework of Ax's generalized Schanuel theorem.
\item \emph{Arithmetic and Differential Galois Groups (2008)}, Oberwolfach Reports, DOI 10.4171/owr/2007/26, pp. 1443--1520. \url{https://doi.org/10.4171/owr/2007/26}.
\item The supplied \texttt{corrected\_verified} note reports editorial and LaTeX cleanup without added source-dependent content. The open label and assessment checked 2026-08-21 say Ax's theorem settles the isoconstant algebraic-group case, while the proposed non-isoconstant extension and differential-Galois interpretation remain open. This is dated supplied triage, not independent verification of its scope or current status.
\end{itemize}
\end{document}
